{\color{red}
\section{Supplementary Materials}\label{sec:reproducibility}\justifying{}
The following supporting information can be downloaded at: https://www.mdpi.com/article/10.3390/pharmaceutics18030287/s1, Figure S1: PRISMA flow diagram; Figure S2: SBML graph of the intestine model of dapagliflozin; Figure S3: SBML graph of the liver model of dapagliflozin; Figure S4: SBML graph of the kidney model of dapagliflozin; Figure S5: Optimization Performance; Figure S6: Simulation Boulton2013; Figure S7: Simulation Cho2021; Figure S8: Simulation FDAMB102006; Figure S9: Simulation Hwang2022a; Figure S10: Simulation Imamura2013; Figure S11: Simulation Jang2020; Figure S12: Simulation Kasichayanula2011c; Figure S13: Simulation Kasichayanula2012; Figure S14: Simulation Kasichayanula2013a; Figure S15: Simulation Khomitskaya2018; Figure S16: Simulation Kim2023; Figure S17: Simulation Kim2023a; Figure S18: Simulation Obermeier2010; Figure S19: Simulation Sha2015; Figure S20: Simulation vanderAartvanderBeek2020; Figure S21: Effect of plasma glucose on pharmacokinetics and pharmacodynamics of dapagliflozin; Figure S22: Local sensitivity analysis; Table S1: Summary of published computational models for dapagliflozin; Table S2: Optimized parameters for the dapagliflozin pharmacokinetic model; Table S3: Optimized parameters for the dapagliflozin pharmacodynamic model.
\section{Author Contributions}
Conceptualization, N.N. and M.K.; methodology, N.N. and M.K.; software, N.N., M.E. and M.K.; validation, N.N., M.E. and M.K.; formal analysis, N.N., M.E. and M.K.; investigation, N.N. and M.K.; resources, M.K.; data curation, N.N., M.E. and M.K.; writing—original draft preparation, N.N. and M.E.; writing—review and editing, N.N., M.E. and M.K.; visualization, N.N., M.E. and M.K.; supervision, M.K.; project administration, M.K.; funding acquisition, M.K. All authors have read and agreed to the published version of the manuscript.
\section{Funding}
Matthias König (MK) was supported by the Federal Ministry of Research, Technology and Space (BMFTR, Germany) within ATLAS by grant number 031L0304B and by the German Research Foundation (DFG) within the Research Unit Program FOR 5151 “QuaLiPerF (Quantifying Liver Perfusion-Function Relationship in Complex Resection—A Systems Medicine Approach)” by grant number 436883643 and by grant number 465194077 (Priority Programme SPP 2311, Subproject SimLivA). This work was supported by the BMBF-funded de.NBI Cloud within the German Network for Bioinformatics Infrastructure (de.NBI) (031A537B, 031A533A, 031A538A, 031A533B, 031A535A, 031A537C, 031A534A, 031A532B).
\section{Institutional Review Board Statement}
Not applicable.
\section{Informed Consent Statement}
Not applicable.
\section{Data Availability Statement}
All curated pharmacokinetic and pharmacodynamic data are publicly available in the PK-DB database (https://pk-db.com) with unique study identifiers as referenced in the manuscript (Table 1).
\section{Acknowledgments}
Figures were created in BioRender. König, M. (2025) https://BioRender.com/07h7x81, accessed on 17 February 2026.
\section{Conflicts of Interest}
The authors declare no conflicts of interest.}